\section{A Linnik obstruction to exclusive prime packets}
\label{sec:prime-packet-linnik-obstruction}

The packet-boundary gate of the preceding section has a second, genuinely
different obstruction.  Its exclusive ownership rule prevents one external
prime from recording simultaneous congruences at several powerful endpoint
primes.  Linnik's least-prime theorem turns this defect into an unconditional
infinite counterfamily.

For a finite packet problem, let the source and sink weights be $x_i\ge0$ and
$y_j\ge0$, and put $X=\sum_i x_i$, $Y=\sum_jy_j$.  If $M_i$ is the mass
assigned to source $i$, then
\[
 \sum_i(x_i-M_i)_+
   =X-\sum_i\min\{x_i,M_i\}.
 \tag{LPO-1}
\]
This is the clipped-reward identity.  In particular, if there are $k$
selected sources with weights $z_1\ge\cdots\ge z_k$ and only $t$ positive
sinks, disjoint ownership lets at most $t$ of those sources receive positive
mass.  Every assignment therefore has residual at least
\[
                         \sum_{i=t+1}^{k}z_i.
 \tag{LPO-2}
\]
For one sink of weight $y$, the exact optimum is
\[
 R_*=X-\max_i\min\{x_i,y\}.
 \tag{LPO-3}
\]
These formulas isolate a cardinality loss which is absent from divisible
transport.

We use Linnik's unconditional theorem in the following uniform form: there
are absolute constants $C\ge1$ and $L\ge1$ such that the least prime in any
reduced residue class $a\pmod Q$ is at most $CQ^L$.  This is precisely the
uniform polynomial estimate proved by Linnik and sharpened, with an explicit
admissible exponent, by Xylouris
\cite[Theorem~1.1]{XylourisLinnik2009}.  No optimal value of $L$ is needed.

Let
\[
 p_1<p_2<\cdots<p_k,\qquad M_k=\prod_{i=1}^{k}p_i,
 \qquad Q_k=M_k^2,
\]
where the $p_i$ are the first $k$ primes.  Apply Linnik's theorem to the
reduced class $-1\pmod{Q_k}$ and choose a prime $\ell_k$ with
\[
 \ell_k\equiv-1\pmod{M_k^2},\qquad
 \ell_k\le C M_k^{2L}.
 \tag{LPO-4}
\]
Then
\[
                         P_k=(1,\ell_k,\ell_k+1)
 \tag{LPO-5}
\]
is a positive primitive $abc$ point.  The family is infinite because
$\ell_k+1\ge M_k^2$.

\begin{lemma}[One sink and many compulsory sources]
\label{lem:linnik-packet-residual}
Every packet assignment at $P_k$ has residual
\[
                         R(P_k)\ge\log M_k-\log p_k.
 \tag{LPO-6}
\]
\end{lemma}

\begin{proof}
The external product is the prime $\ell_k$, so there is exactly one sink.
For every $i\le k$, (LPO-4) gives $p_i^2\mid\ell_k+1$; hence the aggregate
source at $p_i$ has weight
\[
 (v_{p_i}(\ell_k+1)-1)\log p_i\ge\log p_i.
\]
The unique sink can belong to at most one source packet.  Apply (LPO-2) and
discard at most the largest of the $k$ guaranteed weights.
\end{proof}

The conductor remains comparable with $\log M_k$.  Indeed,
\[
 \begin{aligned}
 r(P_k)
 &\le \log\bigl(\ell_k(\ell_k+1)\bigr)\\
 &\le 2\log\ell_k+\log2
 \le4L\log M_k+B,
 \end{aligned}
 \tag{LPO-7}
\]
where $B=2\log C+\log2$.  At the same time,
\[
                         \frac{\log p_k}{\log M_k}\longrightarrow0.
 \tag{LPO-8}
\]
For completeness, Bertrand's postulate gives $p_k<2^k$ for large $k$, while
$p_i\ge i+1$ gives $M_k\ge(k+1)!$; the last half of the factorial factors
already proves (LPO-8).

\begin{theorem}[Exclusive packet obstruction]
\label{thm:linnik-packet-obstruction}
The gate \textup{PBT} is false.
\end{theorem}

\begin{proof}
Set $\eps_0=1/(8L)$.  Combining (LPO-6) and (LPO-7) gives
\[
 R_*(P_k)-\eps_0r(P_k)
 \ge \frac12\log M_k-\log p_k-\frac{B}{8L}.
 \tag{LPO-9}
\]
By (LPO-8), the right side tends to infinity.  Thus every proposed additive
constant fails at some $P_k$, even after minimizing over all assignments.
This has exactly the negated epsilon--constant--point--assignment quantifier
order of PBT.
\end{proof}

This is not an $abc$ counterexample.  Since $M_k$ is even,
$\rad(\ell_k+1)\ge2$, and therefore
\[
 \rad\bigl(\ell_k(\ell_k+1)\bigr)
 =\ell_k\rad(\ell_k+1)\ge2\ell_k\ge\ell_k+1.
 \tag{LPO-10}
\]
Hence every point in the family satisfies $h(P_k)\le r(P_k)$ and has zero
positive scalar defect.  PBT fails solely through its indivisibility gap.
Indeed, if $\Delta(P)=(h(P)-r(P))_+$ and
$G(P)=R_*(P)-\Delta(P)$, then the exact identity
$R_*=\Delta+G$ shows that PBT is equivalent to standard logarithmic $abc$
together with a uniform sublinear bound on $G$.  The Linnik family refutes the
second conjunct while satisfying the first with constant zero.

The correct structural lesson is that one external prime may encode many CRT
incidences at once:
\[
                         \ell_k\equiv-1\pmod{p_i^2}
                         \qquad(1\le i\le k).
\]
A successor must permit this shared incidence while charging its logarithmic
capacity only once.  Allowing arbitrary divisible sharing without incidence
data returns to the scalar defect, so neither repair alone is sufficient.
The next positive target is a shared CRT-incidence boundary with a pointwise
height bridge and an audit showing zero artificial loss on (LPO-5).

The ordinary proof above uses the published Linnik theorem.  The current
Mathlib tree does not contain that analytic theorem.  The companion Lean
module therefore checks the finite packet identities, the prime-neighbour
arithmetic core, and the implication from a transparent Linnik-style escape
hypothesis to the negation of PBT; it does not insert Linnik's theorem as a
project axiom or mislabel the resulting source boundary as kernel-closed.

The accompanying exact computation independently enumerates all
$1{,}368{,}094$ normalized primitive triples with $c\le3000$.  It finds $572$
strict packet fragmentation gaps, $567$ of them with zero scalar defect, and
replays all optima with a second state-set algorithm.  Five finite
square-primorial prime-predecessor rows exhibit the same one-sink mechanism.
These data are regression and discovery evidence only; Theorem~\ref{thm:linnik-packet-obstruction}
uses the infinite Linnik family and does not extrapolate from the scan.
